% ID: SEL-D0078-OB-E02
% Tanda S078.
\textbf{a)} El punto de tangencia es $(1,f(1))=(1,2)$. Derivando, $f'(x)=-2/x^2$, de donde la pendiente es $f'(1)=-2$. La ecuación punto-pendiente da
\[
 y-2=-2(x-1),\qquad\boxed{y=-2x+4}.
\]
\textbf{b)} Para que $(2,5)$ pertenezca a la gráfica, $g(2)=8+4a+b=5$. Al tratarse de un polinomio, un punto de inflexión requiere que se anule la segunda derivada:
\[
 g''(x)=6x+2a,\qquad g''(2)=12+2a=0.
\]
Así, $a=-6$ y, sustituyendo en la primera ecuación, $b=21$. Con estos valores, $g''(x)=6(x-2)$ cambia de negativa a positiva en $2$, confirmando la inflexión. Por tanto, $\boxed{a=-6,\ b=21}$.
